% Part: computability
% Chapter: machines-computations
% Section: configuration

\documentclass[../../../include/open-logic-section]{subfiles}

\begin{document}

\olfileid{cmp}{tur}{con}
\olsection{تشکيلونه او محاسبې}

\begin{explain}
هغه مخکښېنی بحث را ياد کړئ چې د جفت ماشين تشکيلونه مو پکښې يو په بل
پسې څارل۔ د پټې، لوست-ليک سر او ټيورينګ ماشين پروګرام خيالي ميکانيزم
په حقيقت کښې يوازې د ټيورينګ ماشين د محاسبې د ليدلو يوه شهودي لار ده۔
په صوري ډول، پر يوه ټاکلي ننوت د ټيورينګ ماشين محاسبه د
\emph{تشکيلونو} د يوې لړۍ په توګه تعريفولے شو---او تشکيل په خپل وار د
نښو يوه لړۍ (چې د محاسبې په يوه ټاکلي پړاو کښې د پټې له منځپانګې سره
سمون لري)، د لوست-ليک سر د ځاے ښودونکے يو عدد، او يو حالت دے۔ د دې
توکو په وسيله ټاکلے شو چې ټيورينګ ماشين~$M$ پر يوه ټاکلي ننوت څه
محاسبه کوي۔
\end{explain}

\begin{defn}[تشکيل]
د ټيورينګ ماشين $M = \tuple{Q, \Sigma, q_0,
\delta}$ يو \emph{تشکيل} هغه درې‌يز $\tuple{C, m, q}$ دے چې
\begin{enumerate}
\item $C \in \Sigma^*$ د~$\Sigma$ د نښو يوه متناهي لړۍ ده،
\item $m \in \Nat$ يو داسې عدد دے چې $< \len{C}$، او
\item $q \in Q$۔
\end{enumerate}
په شهودي ډول، لړۍ~$C$ د پټې منځپانګه ده (له تر ټولو کيڼ مربع څخه تر
وروستي ناتش يا مخکښې ليدل شوي مربع پورې د ټولو مربعونو نښې)، $m$ د
هغه مربع شمېره ده چې لوست-ليک سر ئې لولي (تر ټولو کيڼ مربع په~$0$
شمېرلو سره)، او~$q$ د ماشين اوسنی حالت دے۔
\end{defn}

\begin{explain}
د ټيورينګ ماشين ممکن ننوت د نښو يوه لړۍ ده، او عموماً داسې لړۍ وي چې
يو عدد په يوه بڼه کوډوي۔ د ټيورينګ ماشين لومړنی تشکيل هغه تشکيل دے
چې ماشين پکښې پر همدغه ننوت کار پيلوي: پټه د پټې د پای نښه لري او
سملاسي ورپسې ښي لوري ته پر مربعونو ننوت ليکل شوے وي، لوست-ليک سر د
ننوت تر ټولو کيڼ مربع (يعنې د کيڼ پای د نښې ښي لوري ته مربع) لولي،
او ميکانيزم په ټاکلي لومړني حالت~$q_0$ کښې وي۔
\end{explain}

\begin{defn}[لومړنی تشکيل]
د ناتش ننوت $I \in \Sigma^*$ دپاره د~$M$ \emph{لومړنی تشکيل} دا دے:
\[
\tuple{\TMendtape \frown I, 1, q_0}.
\]
که ننوت تشه لړۍ وي، لومړنی تشکيل دا دے:
\[
\tuple{\TMendtape \frown \TMblank, 1, q_0}.
\]
\textbf{د ژباړې د نسخې سمون (OLCMP-059):} سرچينه د هرې ممکنې
ننوت-لړۍ دپاره لومړنی تشکيل يوازې د پای نښې او ننوت په يوځاے کولو
جوړوي؛ د تشې ننوت په حالت کښې دې سره د سر شمېره د لړۍ له اوږدوالي
سره برابره شي او د تشکيل شرط ماتېږي۔ دلته د لوستل کېدونکې تشې مربع
په ور زياتولو تشکيل بشپړ شوے دے۔
\end{defn}

\begin{explain}
نښه~$\frown$ د \emph{يوځاے کولو} دپاره ده---د ننوت رشته د کيڼ پای د
نښې سملاسي ښي لوري ته پيلېږي۔
\textbf{د ژباړې د نسخې سمون (OLCMP-045):} د سرچينې عبارت دلته نحوي
نيمګړے دے او د ننوت او د کيڼ پای نښې تر منځ سم لوريز تړاو نۀ
څرګندوي؛ هم پورته نثر او هم صوري لومړنی تشکيل ننوت د نښې ښي لوري ته
ږدي، نو دلته همدغه ټاکلے لوری راوړل شوے دے۔
\end{explain}

\begin{defn}
وايو چې تشکيل $\tuple{C, m, q}$ \emph{په يوه ګام کښې تشکيل
$\tuple{C', m', q'}$ ورکوي} (د~$M$ له مخې)، هله او يوازې هله چې
\begin{enumerate}
\item د~$C$ $m$-مه نښه~$\sigma$ وي،
\item د~$M$ د لارښوونو سټ وټاکي چې $\delta(q, \sigma) =
  \tuple{q', \sigma', D}$،
\item د~$C'$ $m$-مه نښه~$\sigma'$ وي، او
\item
\begin{enumerate}
\item که $m>0$ وي، $D = L$ او $m' = m - 1$ وي، او که نۀ وي $m'=0$، يا
\item $D = R$ او $m' = m + 1$ وي، يا
\item $D = N$ او $m' = m$ وي،
\end{enumerate}
\item که $m' = \len{C}$ وي، نو $\len{C'} = \len{C} + 1$ او د~$C'$
  $m'$-مه نښه~$\TMblank$ ده۔ که نۀ وي، $\len{C'}=\len{C}$۔
\item د هر داسې~$i$ دپاره چې $i < \len{C}$ او $i \neq m$ وي، $C'(i) = C(i)$،
\end{enumerate}
\end{defn}

\begin{defn}\ollabel{defn:run-output}
پر ننوت~$I$ د~$M$ يوه \emph{چلونه} د~$M$ د تشکيلونو داسې لړۍ~$C_i$
ده چې~$C_0$ د ننوت~$I$ دپاره د~$M$ لومړنی تشکيل وي، او هر~$C_i$ په
يوه ګام کښې~$C_{i+1}$ ورکوي۔

وايو چې~$M$ \emph{پر ننوت~$I$ له~$k$ ګامونو وروسته درېږي} که $C_k =
\tuple{C, m, q}$ وي، د~$C$ $m$-مه نښه~$\sigma$ وي، او $\delta(q,
\sigma)$ تعريف شوې نۀ وي۔ په دې صورت کښې د ننوت~$I$ دپاره د~$M$
\emph{وت}~$O$ هله ټاکل کېږي چې~$O$ د نښو داسې رشته وي چې
په~$\TMblank$ نۀ ختمېږي او د کوم~$j \in \Nat$ دپاره $C =
\TMendtape \concat O \concat \TMblank^j$ وي۔ که داسې رشته او طبيعي
عدد شتون ونۀ لري، دې درېدو ته وت نۀ ټاکل کېږي۔ ($\TMblank^j$ د~$j$
شمېر~$\TMblank$ نښو يوه لړۍ ده۔)
\textbf{د ژباړې د نسخې سمون (OLCMP-060):} سرچينه هرې درېدلې چلونې
ته وت ټاکي، که څه هم مخکښې د پټې د پای نښې بدلول روا شوي دي؛ داسې
بدلون ښايي د وت د اړينې بڼې شتون ناممکن کړي۔ دلته وت يوازې هغه وخت
ټاکل کېږي چې دغه بڼه په رښتيا شتون ولري۔
\end{defn}

\begin{explain}
د دې تعريف له مخې د~$M$ هر ټاکل شوی وت~$O$ تل له~$\TMblank$ پرته په
بلې نښې ختمېږي؛ يا که په وخت~$k$ کښې ټوله پټه له~$\TMblank$ څخه ډکه
وي (له تر ټولو کيڼې~$\TMendtape$ پرته)، نو~$O$ تشه رشته ده۔ هغه
درېدلې چلونه چې د پای نښې له اړينې بڼې سره نۀ وي، هېڅ وت نۀ لري۔
\end{explain}

\end{document}
